% Reusable question and answer headings. The second argument is the short
% description displayed in the Table of Contents.
\NewDocumentCommand{\drquestion}{m m +m}{%
  \subsubsection*{\drid{#1}\enspace #3}%
  \addcontentsline{toc}{subsubsection}{#1: #2}%
}
\NewDocumentCommand{\dranswer}{}{%
  \par\noindent{\color{drred}\bfseries Answer:}\par
}

% The shared listings "SQL" language only recognises "--" as a line comment
% and lexes strings by counting quote characters. Our exploited statements
% also use MySQL's "#" line comment (e.g. Q1.1, Q1.3), and because the
% injected payload deliberately leaves an odd, mismatched number of quotes in
% the raw (pre-comment) statement, the automatic lexer cannot reliably tell
% where a "#" comment starts -- it keeps counting quotes as if everything
% were still inside a string, so that text is wrongly coloured as a string
% instead of a comment. (A plain \color escape does not fix this: listings
% re-applies its own per-character style on every following glyph, so a bare
% colour switch never survives past the next character.) The correct fix is
% moredelim's "is" (inline style) delimiter type, which properly integrates
% with listings' own per-character styling engine. Text between a pair of
% "~" characters in ANY lstlisting is forced into the same muted/italic look
% as a native "--" comment, and the "~" characters themselves are not
% printed. Wrap the discarded suffix of a raw statement in "~...~" in every
% future listing that comments via "#" instead of "--" (choose a different,
% unused delimiter character for that one listing if "~" ever needs to
% appear literally in the shown text).
\lstset{moredelim=[is][\color{drmuted}\itshape]{~}{~}}

